\section{The phonon-assisted band sum}\label{app:D.35}

Two one-body vertices acting on a Slater determinant pass only through intermediate states with one electron--hole pair, and on the energy shell the paths through a hole in a filled band take exactly the form of the Pauli-blocked electron paths. We prove the restriction to one-pair intermediate states and the band-sum theorem for the hole paths, both of which \cref{sec:13.2} uses. The proof of the band-sum theorem does not use the form of the vertices, so the theorem holds for light (\cref{sec:6.5}) as well as for dark matter.

\subsection*{Intermediate states and fermion signs}

Let $|\Phi_0\rangle$ be a Slater determinant of one-electron states $|a\rangle$ with energies $\varepsilon_a$: the filled set $\mathcal{F}$ (the valence bands) and the empty set $\mathcal{E}$ (the conduction bands). Let
\begin{equation*}
  X = \sum_{ab}\mathsf{x}_{ab}\,c^\dagger_ac_b , \qquad Y = \sum_{ab}\mathsf{y}_{ab}\,c^\dagger_ac_b ,
\end{equation*}
be two one-body operators. For phonon-assisted electron recoil $\mathsf{x}_{ab}$ is the density vertex $\rho$ and $\mathsf{y}_{ab}$ is $N^{-1/2}g$ times the phonon matrix element $(n_{\vect{Q}\nu} + \frac12 \pm \frac12)^{1/2}$; for light $\mathsf{x}_{ab}$ is the velocity vertex. The process absorbs the energy $\omega$ at the vertex $X$ and changes the energy of the bosons by $\pm\omega_{\vect{Q}\nu}$ at the vertex $Y$. The final state is the single pair
\begin{equation*}
  |\Phi_f\rangle = c^\dagger_cc_v|\Phi_0\rangle , \qquad c\in\mathcal{E},\ v\in\mathcal{F},
\end{equation*}
with $E_f - E_0 = \varepsilon_c - \varepsilon_v \pm \omega_{\vect{Q}\nu}$, the boson energy included. The second-order amplitude is
\begin{equation}
  T = \sum_m\left[\frac{\langle\Phi_f|Y|m\rangle\langle m|X|\Phi_0\rangle}{E_0 + \omega - E_m} + \frac{\langle\Phi_f|X|m\rangle\langle m|Y|\Phi_0\rangle}{E_0 - E_m \mp \omega_{\vect{Q}\nu}}\right] ,
  \label{eq:D.35-T}
\end{equation}
with $E_m$ the electronic energy of the intermediate state; in the second ordering the boson has already been emitted or absorbed. \Cref{eq:D.35-T} is \cref{eq:13.2-T} with the boson energy written out.

\emph{Intermediate states.} Acting on $|\Phi_0\rangle$, $c^\dagger_ac_b$ is nonzero only for $b\in\mathcal{F}$, and then either $a = b$, which returns $|\Phi_0\rangle$, or $a\in\mathcal{E}$, which creates the pair $(a;b)$. So
\begin{equation*}
  X|\Phi_0\rangle = \langle X\rangle_0\,|\Phi_0\rangle + \sum_{a\in\mathcal{E},\,b\in\mathcal{F}}\mathsf{x}_{ab}\,c^\dagger_ac_b|\Phi_0\rangle , \qquad \langle X\rangle_0 = \sum_{b\in\mathcal{F}}\mathsf{x}_{bb} ,
\end{equation*}
and the same for $Y$. Besides the pairs $(p;v)$ with $p\in\mathcal{E}$ and $(c;p)$ with $p\in\mathcal{F}$ of the four paths (\cref{sec:13.2}), the intermediate state can be $|\Phi_0\rangle$ or $|\Phi_f\rangle$ itself, reached through the diagonal parts of the vertices.

\emph{Diagonal terms.} For the density, $\langle X\rangle_0$ is the Fourier component $\bar\rho_e(\vect{q})$ of the ground-state density, nonzero only at $\vect{q} = \vect{G}$; for the electron--phonon vertex the diagonal term needs $\vect{Q} = 0$. When $\langle X\rangle_0$ is nonzero, it enters \cref{eq:D.35-T} twice. With $m = \Phi_0$ in the density-first ordering the term is $\mathsf{y}_{cv}\langle X\rangle_0/\omega$. With $m = \Phi_f$ in the phonon-first ordering the term is $\langle\Phi_f|X|\Phi_f\rangle\,\mathsf{y}_{cv}/(E_0 - E_m \mp \omega_{\vect{Q}\nu})$, where $E_m = E_f \mp \omega_{\vect{Q}\nu}$ is the electronic energy of $\Phi_f$, so that the denominator is $E_0 - E_f = -\omega$ and the term is $-(\langle X\rangle_0 + \mathsf{x}_{cc} - \mathsf{x}_{vv})\,\mathsf{y}_{cv}/\omega$ on the shell $\omega = \varepsilon_c - \varepsilon_v \pm \omega_{\vect{Q}\nu}$. The ground-state expectation value cancels between the two terms, as the linked-cluster theorem requires, and what remains, $(\mathsf{x}_{vv} - \mathsf{x}_{cc})\,\mathsf{y}_{cv}/\omega$, is the $p = v$ electron term of the first ordering plus the $p = c$ electron term of the second, which the band sum below already contains. The diagonal part of $Y$ cancels in the same way.

\subsection*{The band-sum theorem}

Collecting the four kinds of intermediate pair with their energies and the factor $-1$ of each hole path (\cref{sec:13.2}) gives
\begin{equation*}
  \begin{aligned}
  T &= \sum_{p\in\mathcal{E}}\left[\frac{\mathsf{y}_{cp}\,\mathsf{x}_{pv}}{\omega - \varepsilon_p + \varepsilon_v} + \frac{\mathsf{x}_{cp}\,\mathsf{y}_{pv}}{\varepsilon_v - \varepsilon_p \mp \omega_{\vect{Q}\nu}}\right] \\
  &\quad - \sum_{p\in\mathcal{F}}\left[\frac{\mathsf{x}_{cp}\,\mathsf{y}_{pv}}{\omega - \varepsilon_c + \varepsilon_p} + \frac{\mathsf{y}_{cp}\,\mathsf{x}_{pv}}{\varepsilon_p - \varepsilon_c \mp \omega_{\vect{Q}\nu}}\right] ,
  \end{aligned}
\end{equation*}
where the first bracket holds the electron paths A and B of \cref{eq:13.2-paths} and the second the hole paths C and D. On the energy shell, $\omega = \varepsilon_c - \varepsilon_v \pm \omega_{\vect{Q}\nu}$, the hole denominators are
\begin{equation*}
  \omega - \varepsilon_c + \varepsilon_p = -(\varepsilon_v - \varepsilon_p \mp \omega_{\vect{Q}\nu}) , \qquad
  \varepsilon_p - \varepsilon_c \mp \omega_{\vect{Q}\nu} = -(\omega - \varepsilon_p + \varepsilon_v) ,
\end{equation*}
the negatives of the electron denominators with the same $p$. The factor $-1$ of the hole paths and the minus sign of the hole denominators cancel, and the filled-band terms take the form of the empty-band terms:
\begin{equation}
  T = \sum_{p\in\mathcal{E}\cup\mathcal{F}}\left[\frac{\mathsf{y}_{cp}\,\mathsf{x}_{pv}}{\omega - \varepsilon_p + \varepsilon_v} + \frac{\mathsf{x}_{cp}\,\mathsf{y}_{pv}}{\varepsilon_v - \varepsilon_p \mp \omega_{\vect{Q}\nu}}\right]
  = \langle c|\,\mathsf{y}\,\frac{1}{\varepsilon_v + \omega - h}\,\mathsf{x} + \mathsf{x}\,\frac{1}{\varepsilon_v \mp \omega_{\vect{Q}\nu} - h}\,\mathsf{y}\,|v\rangle ,
  \label{eq:D.35-theorem}
\end{equation}
where $h$ is the one-electron Hamiltonian and $\mathsf{x}$, $\mathsf{y}$ are the one-electron operators with matrix elements $\mathsf{x}_{ab}$, $\mathsf{y}_{ab}$. The terms $p = v$ in the first resolvent and $p = c$ in the second are finite, with denominators $\omega$ and $-\omega$. With $\mathsf{x} = \rho$, $\mathsf{y} = g$ (the factor $N^{-1/2}(n_{\vect{Q}\nu} + \frac12 \pm \frac12)^{1/2}$ taken out) and the crystal momenta written out, \cref{eq:D.35-theorem} is \cref{eq:13.2-M}; with $\mathsf{x}$ the velocity vertex it is \cref{eq:6.5-M}.

\subsection*{The error of the adiabatic form}

The full amplitude \cref{eq:13.2-M} differs from the adiabatic amplitude, \cref{eq:13.2-Mad}, only by $\pm\omega_{\vect{Q}\nu}$ in the denominators. On the shell the first denominator is $D_{\mathrm{A}} \pm \omega_{\vect{Q}\nu}$ with $D_{\mathrm{A}} = \varepsilon_{c\vect{k}'} - \varepsilon_{p,\vect{k}+\vect{q}}$, and the second $D_{\mathrm{B}} \mp \omega_{\vect{Q}\nu}$ with $D_{\mathrm{B}} = \varepsilon_{v\vect{k}} - \varepsilon_{p,\vect{k}+\vect{Q}}$. With $1/(D \pm \omega) = 1/D \mp \omega/D^2 + \omega^2/D^3 \mp \cdots$,
\begin{equation*}
  M^\pm - M^{\mathrm{ad}} = \mp\,\omega_{\vect{Q}\nu}\sum_p\left[\frac{g_{cp\nu}\,\rho_{pv}}{D_{\mathrm{A}}^2} - \frac{\rho_{cp}\,g_{pv\nu}}{D_{\mathrm{B}}^2}\right] + \omega_{\vect{Q}\nu}^2\sum_p\left[\frac{g_{cp\nu}\,\rho_{pv}}{D_{\mathrm{A}}^3} + \frac{\rho_{cp}\,g_{pv\nu}}{D_{\mathrm{B}}^3}\right] + \cdots ,
\end{equation*}
which is \cref{eq:13.2-err}. Let $\delta_{\mathrm{off}}$ be the smallest of the full denominators $|D_{\mathrm{A}} \pm \omega_{\vect{Q}\nu}|$ and $|D_{\mathrm{B}} \mp \omega_{\vect{Q}\nu}|$, bounded below by \cref{eq:13.2-doff}. Each adiabatic denominator then exceeds $\delta_{\mathrm{off}} - \omega_{\vect{Q}\nu}$, the $n$th term of the series is bounded by $(\omega_{\vect{Q}\nu}/(\delta_{\mathrm{off}} - \omega_{\vect{Q}\nu}))^n$ times $\sum_p(|g_{cp\nu}\rho_{pv}|/|D_{\mathrm{A}}| + |\rho_{cp}g_{pv\nu}|/|D_{\mathrm{B}}|)$, and
\begin{equation*}
  \bigl|M^\pm - M^{\mathrm{ad}}\bigr| \le \frac{\omega_{\vect{Q}\nu}}{\delta_{\mathrm{off}} - 2\omega_{\vect{Q}\nu}}\sum_p\left[\frac{|g_{cp\nu}\,\rho_{pv}|}{|D_{\mathrm{A}}|} + \frac{|\rho_{cp}\,g_{pv\nu}|}{|D_{\mathrm{B}}|}\right] ,
\end{equation*}
for $\delta_{\mathrm{off}} > 2\omega_{\vect{Q}\nu}$. The sum on the right is $|M^{\mathrm{ad}}|$ when its terms do not cancel, and the relative error is then of order $\omega_{\vect{Q}\nu}/\delta_{\mathrm{off}}$. When the terms cancel strongly, as for $\vect{q}\to0$ where $M^\pm$ and $M^{\mathrm{ad}}$ both vanish by \cref{eq:13.2-q0}, the bound is on the absolute error only.

\subsection*{The rate and the loss function}

The operator $B_{\vect{Q}\nu} = b_{\vect{Q}\nu} + b^\dagger_{-\vect{Q}\nu}$ changes one occupation by one, so final states that differ in which mode changed are orthogonal. Averaging the squared matrix elements mode by mode over the thermal state of the lattice gives $n_{\vect{Q}\nu} + 1$ for emission and $n_{\vect{Q}\nu}$ for absorption, $n_{\vect{Q}\nu} = n_{\mathrm{B}}(\omega_{\vect{Q}\nu})$: the factors $n_{\vect{Q}\nu} + \frac12 \pm \frac12$ of \cref{eq:13.2-M}. For the electrons, the transition from $(v,\vect{k})$ to $(c,\vect{k}')$ needs the first state filled and the second empty, with the weight $f_v(1 - f_c)$ of the Fermi occupations; with the chemical potential in the gap, the weight $f_v(1 - f_c)$ differs from one by terms of order $\ee^{-\Eg/2T}$. The squared amplitude goes into \cref{eq:2.5-golden}. Summing over $\vect{k}$, $\vect{Q}$, the bands, the branches, the sign and the two spin states, and using $\Vol = N\Omega$ with \cref{eq:A.3-BZ}, gives \cref{eq:13.3-SPA} with the extra factor $f_v(1 - f_c)$ inside the average.

The loss function is the dissipative part of the response, built from the rate of each transition minus the rate of the inverse transition. For emission, the inverse transition removes the pair and absorbs a phonon, with the weight $f_c(1 - f_v)\,n_{\vect{Q}\nu}$. In equilibrium $f_c(1 - f_v) = f_v(1 - f_c)\,\ee^{-\beta(\varepsilon_{c\vect{k}'} - \varepsilon_{v\vect{k}})}$ and $n_{\vect{Q}\nu} = (n_{\vect{Q}\nu} + 1)\,\ee^{-\beta\omega_{\vect{Q}\nu}}$, so on the energy shell
\begin{equation*}
  f_v(1 - f_c)(n_{\vect{Q}\nu} + 1) - f_c(1 - f_v)\,n_{\vect{Q}\nu} = f_v(1 - f_c)(n_{\vect{Q}\nu} + 1)\bigl(1 - \ee^{-\beta\omega}\bigr) ,
\end{equation*}
and in the same way for absorption with $n_{\vect{Q}\nu}$ in place of $n_{\vect{Q}\nu} + 1$. The factor $1 - \ee^{-\beta\omega}$ of \cref{eq:4.4-SL} is therefore carried by each transition, and $L^{\mathrm{PA}}$ is \cref{eq:13.3-LPA} with the weights $w_\pm f_v(1 - f_c)$. Both weights are nonnegative for $\omega > 0$, and so is $L^{\mathrm{PA}}$. For $T \ll \Eg$ the Fermi factors are one.
